% GENERATED FILE. Edit canonical lessons, then run npm run generate:books.
% canonical-source-hash: fnv1a64:954a639b66228be6
% canonical-lessons: SA-C104-satika, SA-C104-padatranam, SA-C104-usnisam, SA-C104-kancukah, SA-C104-durabhasah

\chapter{Sari, Shoe, Turban, Shirt, Telephone}
\label{ch:sa-sari-shoe-turban-shirt-telephone}
\hlchaptermodality{104}

\begin{chapteropening}
\textbf{By the end of this chapter:} \emph{I can say a sari, a shoe, a turban, a shirt and a telephone.}

Complete the last lesson of chapter 104: i can say a sari, a shoe, a turban, a shirt and a telephone.
\end{chapteropening}

\section[\texorpdfstring{\sk{शाटिका} (śāṭikā)}{शाटिका (śāṭikā)}]{\sk{शाटिका} (śāṭikā) — a sari}

\label{lesson:SA-C104-satika}



\begin{quote}
Before the new one: say the Sanskrit for a bangle, then the Sanskrit for a shawl.
\end{quote}

\subsection*{You'll want to know: \sk{शाटिका}}
\textbf{\sk{शाटिका}} — \emph{śāṭikā} — \textquotedblleft{}a sari\textquotedblright{}.

\begin{grammarlens}[title={What you've built}]
One more word for this chapter.
\end{grammarlens}

\subsection*{Your turn}
\begin{itemize}
\raggedright
  \item \emph{Say it:} \emph{śāṭikā}
  \item \emph{Say it:} \emph{śāṭikā}, once more
  \item \emph{Say it:} say \emph{śāṭikā}
  \item \emph{Recall:} say the Sanskrit for a bangle, then the Sanskrit for a shawl, then say \emph{śāṭikā} again
\end{itemize}

\begin{tcolorbox}[breakable,colback=teal!4,colframe=teal!35!black,title={Before you move on}]
What does \sk{शाटिका} mean? (\textquotedblleft{}A sari\textquotedblright{}.) Say it once more.
\end{tcolorbox}
\section[\texorpdfstring{\sk{पादत्राणम्} (pādatrāṇam)}{पादत्राणम् (pādatrāṇam)}]{\sk{पादत्राणम्} (pādatrāṇam) — a shoe}

\label{lesson:SA-C104-padatranam}



\begin{quote}
Before the new one: say the Sanskrit for a shawl, then the Sanskrit for a sari.
\end{quote}

\subsection*{You'll want to know: \sk{पादत्राणम्}}
\textbf{\sk{पादत्राणम्}} — \emph{pādatrāṇam} — \textquotedblleft{}a shoe\textquotedblright{}.

\begin{grammarlens}[title={What you've built}]
One more word for this chapter.
\end{grammarlens}

\subsection*{Your turn}
\begin{itemize}
\raggedright
  \item \emph{Say it:} \emph{pādatrāṇam}
  \item \emph{Say it:} \emph{pādatrāṇam}, once more
  \item \emph{Say it:} say \emph{pādatrāṇam}
  \item \emph{Recall:} say the Sanskrit for a shawl, then the Sanskrit for a sari, then say \emph{pādatrāṇam} again
\end{itemize}

\begin{tcolorbox}[breakable,colback=teal!4,colframe=teal!35!black,title={Before you move on}]
What does \sk{पादत्राणम्} mean? (\textquotedblleft{}A shoe\textquotedblright{}.) Say it once more.
\end{tcolorbox}
\section[\texorpdfstring{\sk{उष्णीषम्} (uṣṇīṣam)}{उष्णीषम् (uṣṇīṣam)}]{\sk{उष्णीषम्} (uṣṇīṣam) — a turban}

\label{lesson:SA-C104-usnisam}



\begin{quote}
Before the new one: say the Sanskrit for a sari, then the Sanskrit for a shoe.
\end{quote}

\subsection*{You'll want to know: \sk{उष्णीषम्}}
\textbf{\sk{उष्णीषम्}} — \emph{uṣṇīṣam} — \textquotedblleft{}a turban\textquotedblright{}.

\begin{grammarlens}[title={What you've built}]
One more word for this chapter.
\end{grammarlens}

\subsection*{Your turn}
\begin{itemize}
\raggedright
  \item \emph{Say it:} \emph{uṣṇīṣam}
  \item \emph{Say it:} \emph{uṣṇīṣam}, once more
  \item \emph{Say it:} say \emph{uṣṇīṣam}
  \item \emph{Recall:} say the Sanskrit for a sari, then the Sanskrit for a shoe, then say \emph{uṣṇīṣam} again
\end{itemize}

\begin{tcolorbox}[breakable,colback=teal!4,colframe=teal!35!black,title={Before you move on}]
What does \sk{उष्णीषम्} mean? (\textquotedblleft{}A turban\textquotedblright{}.) Say it once more.
\end{tcolorbox}
\section[\texorpdfstring{\sk{कञ्चुकः} (kañcukaḥ)}{कञ्चुकः (kañcukaḥ)}]{\sk{कञ्चुकः} (kañcukaḥ) — a shirt, a blouse}

\label{lesson:SA-C104-kancukah}



\begin{quote}
Before the new one: say the Sanskrit for a shoe, then the Sanskrit for a turban.
\end{quote}

\subsection*{You'll want to know: \sk{कञ्चुकः}}
\textbf{\sk{कञ्चुकः}} — \emph{kañcukaḥ} — \textquotedblleft{}a shirt, a blouse\textquotedblright{}.

\begin{grammarlens}[title={What you've built}]
One more word for this chapter.
\end{grammarlens}

\subsection*{Your turn}
\begin{itemize}
\raggedright
  \item \emph{Say it:} \emph{kañcukaḥ}
  \item \emph{Say it:} \emph{kañcukaḥ}, once more
  \item \emph{Say it:} say \emph{kañcukaḥ}
  \item \emph{Recall:} say the Sanskrit for a shoe, then the Sanskrit for a turban, then say \emph{kañcukaḥ} again
\end{itemize}

\begin{tcolorbox}[breakable,colback=teal!4,colframe=teal!35!black,title={Before you move on}]
What does \sk{कञ्चुकः} mean? (\textquotedblleft{}A shirt, a blouse\textquotedblright{}.) Say it once more.
\end{tcolorbox}
\section[\texorpdfstring{\sk{दूरभाषः} (dūrabhāṣaḥ)}{दूरभाषः (dūrabhāṣaḥ)}]{\sk{दूरभाषः} (dūrabhāṣaḥ) — a telephone}

\label{lesson:SA-C104-durabhasah}



\begin{quote}
Before the new one: say the Sanskrit for a turban, then the Sanskrit for a shirt.
\end{quote}

\subsection*{You'll want to know: \sk{दूरभाषः}}
\textbf{\sk{दूरभाषः}} — \emph{dūrabhāṣaḥ} — \textquotedblleft{}a telephone\textquotedblright{}.

\begin{grammarlens}[title={What you've built}]
One more word for this chapter.
\end{grammarlens}

\subsection*{Your turn}
\begin{itemize}
\raggedright
  \item \emph{Say it:} \emph{dūrabhāṣaḥ}
  \item \emph{Say it:} \emph{dūrabhāṣaḥ}, once more
  \item \emph{Say it:} say \emph{dūrabhāṣaḥ}
  \item \emph{Recall:} say the Sanskrit for a turban, then the Sanskrit for a shirt, then say \emph{dūrabhāṣaḥ} again
\end{itemize}

\begin{tcolorbox}[breakable,colback=teal!4,colframe=teal!35!black,title={Before you move on}]
What does \sk{दूरभाषः} mean? (\textquotedblleft{}A telephone\textquotedblright{}.) Say it once more.
\end{tcolorbox}
